
\documentclass{article}
\usepackage{colortbl}
\usepackage{makecell}
\usepackage{multirow}
\usepackage{supertabular}

\begin{document}

\newcounter{utterance}

\centering \large Interaction Transcript for game `cladder', experiment `full\_v1.5\_default', episode 2287 with Qwen3.5{-}9B.
\vspace{24pt}

{ \footnotesize  \setcounter{utterance}{1}
\setlength{\tabcolsep}{0pt}
\begin{supertabular}{c@{$\;$}|p{.15\linewidth}@{}p{.15\linewidth}p{.15\linewidth}p{.15\linewidth}p{.15\linewidth}p{.15\linewidth}}
    \# & \multicolumn{2}{c}{Player} && \multicolumn{2}{c}{Game Master} \\
    \hline

    \theutterance \stepcounter{utterance}  
    & & & \multicolumn{4}{p{0.6\linewidth}}{
        \cellcolor[rgb]{0.9,0.9,0.9}{
            \makecell[{{p{\linewidth}}}]{
                \texttt{\tiny{[P1$\langle$GM]}}
                \texttt{You are an expert in causal inference. The following question is not a typical commonsense query, but rather a meticulously designed question created by a professor specializing in causal inference, intended to assess the students' mastery of the course content.} \\
\texttt{Imagine a self{-}contained, hypothetical world with only the following conditions, and without any unmentioned factors or causal relationships: Smoking has a direct effect on effort and college admission. Effort has a direct effect on college admission. The overall probability of smoking is 93\%. The probability of nonsmoking and college admission is 3\%. The probability of smoking and college admission is 62\%. Is the chance of college admission smaller when observing smoking?} \\
\texttt{Start your answer with 'Yes' or 'No', followed by additional reasoning or evidence to support your explanation.} \\
            }
        }
    }
     \\ \\

    \theutterance \stepcounter{utterance}  
    & \multicolumn{4}{p{0.6\linewidth}}{
        \cellcolor[rgb]{0.561,0.737,0.561}{
            \makecell[{{p{\linewidth}}}]{
                \texttt{\tiny{[P1$\rangle$GM]}}
                \texttt{No} \\
\\ 
\texttt{**Reasoning:**} \\
\\ 
\texttt{To determine whether the chance of college admission is smaller when observing smoking, we must compare the conditional probability of admission given smoking, \$P(\textbackslash{}text\{Admission\} | \textbackslash{}text\{Smoking\})\$, with the unconditional probability of admission, \$P(\textbackslash{}text\{Admission\})\$.} \\
\\ 
\texttt{Based on the provided data:} \\
\texttt{1.  **Calculate the unconditional probability of admission, \$P(\textbackslash{}text\{Admission\})\$:**} \\
\texttt{    The problem states that the probability of nonsmoking and college admission is 3\% (\$P(\textbackslash{}neg S \textbackslash{}cap A) = 0.03\$) and the probability of smoking and college admission is 62\% (\$P(S \textbackslash{}cap A) = 0.62\$).} \\
\texttt{    Since these are the only two mutually exclusive ways to be admitted (smoking or not smoking), the total probability is the sum:} \\
\texttt{    \$\$P(\textbackslash{}text\{Admission\}) = P(S \textbackslash{}cap A) + P(\textbackslash{}neg S \textbackslash{}cap A) = 0.62 + 0.03 = 0.65\$\$} \\
\texttt{    So, \$P(\textbackslash{}text\{Admission\}) = 65\textbackslash{}\%\$.} \\
\\ 
\texttt{2.  **Calculate the conditional probability of admission given smoking, \$P(\textbackslash{}text\{Admission\} | \textbackslash{}text\{Smoking\})\$:**} \\
\texttt{    We are given that the overall probability of smoking is 93\% (\$P(S) = 0.93\$).} \\
\texttt{    Using the definition of conditional probability:} \\
\texttt{    \$\$P(\textbackslash{}text\{Admission\} | \textbackslash{}text\{Smoking\}) = \textbackslash{}frac\{P(\textbackslash{}text\{Smoking\} \textbackslash{}cap \textbackslash{}text\{Admission\})\}\{P(\textbackslash{}text\{Smoking\})\}\$\$} \\
\texttt{    \$\$P(\textbackslash{}text\{Admission\} | \textbackslash{}text\{Smoking\}) = \textbackslash{}frac\{0.62\}\{0.93\} \textbackslash{}approx 0.6667\$\$} \\
\texttt{    So, \$P(\textbackslash{}text\{Admission\} | \textbackslash{}text\{Smoking\}) \textbackslash{}approx 66.7\textbackslash{}\%\$.} \\
\\ 
\texttt{3.  **Comparison:**} \\
\texttt{    The question asks if the chance of admission is *smaller* when observing smoking.} \\
\texttt{    Comparing the two values: \$66.7\textbackslash{}\% > 65\textbackslash{}\%\$.} \\
\texttt{    The probability of admission increases (from 65\% to \textasciitilde{}66.7\%) when we observe that an individual smokes.} \\
\\ 
\texttt{Therefore, the chance of college admission is **not** smaller; it is slightly larger.} \\
            }
        }
    }
    & & \\ \\

    \theutterance \stepcounter{utterance}  
    & & & \multicolumn{2}{p{0.3\linewidth}}{
        \cellcolor[rgb]{0.9,0.9,0.9}{
            \makecell[{{p{\linewidth}}}]{
                \texttt{\tiny{[GM$|$GM]}}
                \texttt{No} \\
            }
        }
    }
    & & \\ \\

    \theutterance \stepcounter{utterance}  
    & & & \multicolumn{2}{p{0.3\linewidth}}{
        \cellcolor[rgb]{0.9,0.9,0.9}{
            \makecell[{{p{\linewidth}}}]{
                \texttt{\tiny{[GM$|$GM]}}
                \texttt{no} \\
            }
        }
    }
    & & \\ \\

    \theutterance \stepcounter{utterance}  
    & & & \multicolumn{2}{p{0.3\linewidth}}{
        \cellcolor[rgb]{0.9,0.9,0.9}{
            \makecell[{{p{\linewidth}}}]{
                \texttt{\tiny{[GM$|$GM]}}
                \texttt{game\_result = WIN} \\
            }
        }
    }
    & & \\ \\

\end{supertabular}
}

\end{document}
